% from ega0-1-fr.tex
\oldpage[0\textsubscript{I}]{73}

\begin{corollary}[7.6.4]
\label{0.7.6.4-fr}
Si, dans $A$, $0$ n'est pas adhérent à $S$, l'anneau $A\{S^{-1}\}$ n'est pas
réduit à $0$.
\end{corollary}

En effet, $0$ n'est pas adhérent à $\{1\}$ dans l'anneau $S^{-1}A$~; sinon,
on aurait $1\in S^{-1}\mathfrak J_\lambda$ pour tout idéal ouvert
$\mathfrak J_\lambda$ de $A$, et il en résulterait que
$\mathfrak J_\lambda\cap S\neq\varnothing$ pour tout $\lambda$, contrairement
à l'hypothèse.

\begin{env}[7.6.5]
\label{0.7.6.5-fr}
Nous dirons que $A\{S^{-1}\}$ est l'\emph{anneau complet des fractions} de $A$
ayant leurs dénominateurs dans $S$. Avec les notations précédentes, il est clair
que l'image réciproque de $S^{-1}\mathfrak J_\lambda$ dans $A$ contient
$\mathfrak J_\lambda$, donc l'application canonique $A\to S^{-1}A$ est
continue, et si on la compose avec l'application canonique
$S^{-1}A\to A\{S^{-1}\}$, on obtient un homomorphisme canonique continu
$A\to A\{S^{-1}\}$, limite projective des homomorphismes
$A\to S_\lambda^{-1}A_\lambda$.
\end{env}

\begin{env}[7.6.6]
\label{0.7.6.6-fr}
Le couple formé de $A\{S^{-1}\}$ et de l'application canonique
$A\to A\{S^{-1}\}$ est caractérisé par la \emph{propriété universelle}
suivante~: tout homomorphisme continu $u$ de $A$ dans un anneau linéairement
topologisé $B$, séparé et complet, tel que $u(S)$ soit formé d'éléments
inversibles dans $B$, se factorise d'une seule manière en
$A\to A\{S^{-1}\}\xrightarrow{u'}B$, où $u'$ est continu. En effet, $u$ se
factorise d'une seule manière en $A\to S^{-1}A\xrightarrow{v'}B$~; comme pour
tout idéal ouvert $\mathfrak K$ de $B$, $u^{-1}(\mathfrak K)$ contient un
$\mathfrak J_\lambda$, ${v'}^{-1}(\mathfrak K)$ contient nécessairement
$S^{-1}\mathfrak J_\lambda$, donc $v'$ est continu~; puisque $B$ est séparé et
complet, $v'$ se factorise d'une seule manière en
$S^{-1}A\to A\{S^{-1}\}\xrightarrow{u'}B$, où $u'$ est continu~; d'où notre
assertion.
\end{env}

\begin{env}[7.6.7]
\label{0.7.6.7-fr}
Soient $B$ un second anneau linéairement topologisé, $T$ une partie
multiplicative de $B$, $\varphi:A\to B$ un homomorphisme continu tel que
$\varphi(S)\subset T$. D'après ce qui précède, l'homomorphisme continu
$A\xrightarrow{\varphi}B\to B\{T^{-1}\}$ se factorise de façon unique en
$A\to A\{S^{-1}\}\xrightarrow{\varphi'}B\{T^{-1}\}$, où $\varphi'$ est
continu. En particulier, si $B=A$ et si $\varphi$ est l'identité, on voit que
pour $S\subset T$ on a un homomorphisme continu
$\rho^{T,S}:A\{S^{-1}\}\to A\{T^{-1}\}$ obtenu par passage au séparé complété
à partir de $S^{-1}A\to T^{-1}A$~; si $U$ est une troisième partie
multiplicative de $A$ telle que $S\subset T\subset U$, on a
$\rho^{U,S}=\rho^{U,T}\circ\rho^{T,S}$.
\end{env}

\begin{env}[7.6.8]
\label{0.7.6.8-fr}
Soient $S_1$, $S_2$ deux parties multiplicatives de $A$, et soit $S_2'$ l'image
canonique de $S_2$ dans $A\{S_1^{-1}\}$~; on a alors un isomorphisme
topologique canonique
$A\{(S_1S_2)^{-1}\}\xrightarrow{\sim}A\{S_1^{-1}\}\{{S_2'}^{-1}\}$,
comme on le voit en partant de l'isomorphisme canonique
$(S_1S_2)^{-1}A\xrightarrow{\sim}{S_2''}^{-1}(S_1^{-1}A)$ (où $S_2''$ est
l'image canonique de $S_2$ dans $S_1^{-1}A$), qui est bicontinu.
\end{env}

\begin{env}[7.6.9]
\label{0.7.6.9-fr}
Soit $\mathfrak a$ un idéal \emph{ouvert} de $A$~; on peut supposer que
$\mathfrak J_\lambda\subset\mathfrak a$ pour tout $\lambda$, et par suite
$S^{-1}\mathfrak J_\lambda\subset S^{-1}\mathfrak a$ dans l'anneau
$S^{-1}A$, autrement dit, $S^{-1}\mathfrak a$ est un idéal \emph{ouvert} de
$S^{-1}A$~; nous désignerons par $\mathfrak a\{S^{-1}\}$ son séparé complété,
égal à $\varprojlim(S^{-1}\mathfrak a/S^{-1}\mathfrak J_\lambda)$, qui est un
idéal \emph{ouvert} de $A\{S^{-1}\}$, isomorphe à l'adhérence de l'image
canonique de $S^{-1}\mathfrak a$. En outre, \emph{l'anneau discret
$A\{S^{-1}\}/\mathfrak a\{S^{-1}\}$ est canoniquement isomorphe à
$S^{-1}A/S^{-1}\mathfrak a=S^{-1}(A/\mathfrak a)$}. Inversement, si
$\mathfrak a'$ est un idéal ouvert de $A\{S^{-1}\}$, $\mathfrak a'$ contient
un idéal de la forme $\mathfrak J_\lambda\{S^{-1}\}$, donc est l'image
réciproque d'un idéal de $S^{-1}A/S^{-1}\mathfrak J_\lambda$, qui est
nécessairement (1.2.6) de la forme $S^{-1}\mathfrak a$, où
$\mathfrak a\supset\mathfrak J_\lambda$. On en conclut que l'on a
$\mathfrak a'=\mathfrak a\{S^{-1}\}$. En particulier (1.2.6)~:
\end{env}

\begin{proposition}[7.6.10]
\label{0.7.6.10-fr}
L'application $\mathfrak p\mapsto\mathfrak p\{S^{-1}\}$ est une bijection
croissante de l'ensemble des idéaux premiers ouverts $\mathfrak p$ de $A$ tels
que $\mathfrak p\cap S=\varnothing$ sur l'ensemble des idéaux premiers ouverts

